\documentclass[12pt]{article}
\usepackage[margin=1in]{geometry}
\usepackage{setspace}
\usepackage{amsmath, amssymb, amsthm}
\usepackage{hyperref}
\usepackage{enumitem}

\title{Reviewer Report\\ \large \textit{Political Polarization and the Interpretation of Auditor Signals}}
\author{Prepared for the 2026 Conference on Auditing and Capital Markets}
\date{\today}

\begin{document}
\maketitle
\onehalfspacing

\section{General Evaluation}

The manuscript investigates a highly relevant question: does political polarization affect how investors interpret auditor-related information? The premise that polarization increases dispersion in prior beliefs about institutional credibility is intuitive and economically important. However, as currently written, the draft suffers from a mechanical application of polarization theory and an empirical proxy that is outdated. To reach the standard of a top-tier accounting journal, the paper must articulate the precise psychological mechanism driving this dispersion and deploy more robust mathematical and empirical specifications. 

The integration of recent advancements in the measurement of polarization---specifically regarding affective polarization and cluster-based variance---will be critical for the next iteration of this draft.

\section{Major Comments: Theory and Mechanism}

\subsection{Decomposing the Mechanism of Distrust}
The draft argues that polarization reduces consensus about the credibility of auditors. However, it treats polarization as a monolith. Drawing on Campos and Federico (2025), affective polarization is actually a multidimensional construct comprising othering, aversion, and moralization[cite: 319, 394]. 

Your theoretical framework would be significantly strengthened by identifying which of these components drives the heterogeneous interpretation of audit signals. For example, Campos and Federico (2025) note that \textit{aversion} (the tendency to dislike and avoid out-party members) is strongly tied to an ``anti-establishment'' orientation[cite: 393, 782]. Because auditing is a regulated, establishment-dependent institution, an aversion-driven rejection of democratic or establishment norms perfectly explains why highly polarized investors would discount the standardized likelihood function of an audit report[cite: 756, 785]. By contrast, \textit{moralization} often aligns with an affirmation of the current political establishment[cite: 784]. Refining Section 3 to specify that \textit{aversion}-based polarization drives the erosion of institutional trust will give your model the theoretical precision expected by rigorous reviewers.

\subsection{Elevating the Mathematical Rigor of Belief Updating}
The current Bayesian updating model (Equations 1-3) is a bit pedestrian for a top-tier submission. Using a simple binary state space and a discrete signal ($\theta \in \{0, 1\}$, $s \in \{0, 1\}$) feels underdeveloped when mapped to continuous variables like trading volume and absolute cumulative abnormal returns ($|CAR|$). 

I strongly suggest extending this model to a continuous state space. Formalizing the belief distributions using a stochastic process or a functional analysis approach to the prior distributions would elevate the theoretical rigor. You could model the audit signal as a diffusion process where polarization acts as a modifier on the variance of the drift parameter. I highly recommend simulating this continuous belief updating process and the resulting trading volume dynamics in R or Python to validate your comparative statics before taking them to the data.

\section{Major Comments: Measurement and Empirics}

\subsection{Limitations of the Esteban and Ray (1994) Index}
The current empirical design relies heavily on the Esteban and Ray (1994) index applied to U.S. House election returns (Equation 6). Mehlhaff (2023) points out severe limitations with this measure. Specifically, the Esteban and Ray index is based on arbitrary assumptions and is shown to be essentially indistinguishable from simple variance, with correlations as high as 0.933[cite: 228, 229]. 

Furthermore, Mehlhaff (2023) argues that standard measures of polarization fail because they only capture intergroup heterogeneity while ignoring intragroup homogeneity[cite: 8, 14, 221]. If you continue to use voting data, you must address the fact that party sorting (intragroup homogeneity) heavily influences these metrics[cite: 34, 44]. 

\subsection{Adopting or Benchmarking Against the Cluster-Polarization Coefficient (CPC)}
To resolve the measurement issue, you should evaluate your data using the Cluster-Polarization Coefficient (CPC) introduced by Mehlhaff (2023)[cite: 9, 57]. The CPC explicitly models both intergroup heterogeneity (Between-Cluster Sum of Squares) and intragroup homogeneity (Within-Cluster Sum of Squares) by decomposing total variance[cite: 59, 93].

Because you are using congressional district voting aggregated to the state level, adopting the CPC would allow you to accurately capture situations where a state is polarized not just because the vote split is 50/50, but because the voters within those coalitions are highly homogeneous and geographically/ideologically concentrated. The CPC is also unbiased and consistent in small samples due to corrections for lost degrees of freedom[cite: 76, 85], which is crucial if you are analyzing state-by-year panels.

\section{Recommendations for the Next Revision}
\begin{enumerate}[label=\textbf{\arabic*.}]
    \item \textbf{Revise Section 3 (Theory):} Incorporate the multidimensionality of affective polarization (aversion vs. moralization) to explain \textit{why} investors distrust the audit signal[cite: 319, 782]. Upgrade the Bayesian updating model to a continuous distribution.
    \item \textbf{Revise Section 5 (Empirical Design):} Either replace the Esteban and Ray (1994) measure with the CPC [cite: 63], or calculate the CPC as a primary robustness check (your proposed Table 5) to prove your results are not an artifact of simple variance[cite: 228].
    \item \textbf{Simulate the Model:} Build an R or Python script to simulate the theoretical market responses based on heterogeneous priors. Include these simulation plots to validate Hypothesis 3 (Signal Ambiguity).
\end{enumerate}

\end{document}